Comparing two group means is one of the most common analyses in empirical research. The usual tool is Student's $t$-test~\citep{student1908probable}, which pools the two sample variances into one estimate. Pooling is valid when both groups share a variance. When they do not, the test statistic no longer follows the assumed $t$ distribution, and the error rate drifts away from the level the analyst chose.

Welch~\citep{welch1947generalization} proposed a version that estimates each variance separately and adjusts the degrees of freedom. Methodologists have argued for years that Welch's test should replace Student's test as the default~\citep{ruxton2006unequal, delacre2017why}. The argument rests on a specific failure: Student's test becomes too liberal when the smaller group has the larger variance, because the pooled estimate is dominated by the larger, quieter group.

We quantify that failure in a setting that is easy to reproduce. Our contribution is a small, fully scripted simulation that reports both error rates with confidence intervals and a paired test of the difference. The script and its output are part of this paper, so every number below can be regenerated with one command.
